\documentclass[margin, line]{res}
%\usepackage[margin=0.5in]{geometry}
\usepackage{url}
\usepackage[normalem]{ulem}
\usepackage{hyperref}
\usepackage{ctex}
\usepackage{bbding}
%\topmargin=-0.5in
%\bottommargin=-0.5in
\oddsidemargin
-.5in
\evensidemargin
-.5in
\textwidth=5.6in
\textheight=9.1in
\itemsep=0.5in
\parsep=0.5in

\newenvironment{list1}{ \begin{list}{\ding{113}}{%
\setlength{\itemsep}{0in} \setlength{\parsep}{0in} \setlength{\parskip}{0in} \setlength{\topsep}{0in} \setlength{\partopsep}{0in} \setlength{\leftmargin}{0.17in}}}{\end{list}}
\newenvironment{list2}{ \begin{list}{$\bullet$}{%
\setlength{\itemsep}{0in} \setlength{\parsep}{0in} \setlength{\parskip}{0in} \setlength{\topsep}{0in} \setlength{\partopsep}{0in} \setlength{\leftmargin}{0.2in}}}{\end{list}}

\pagestyle{plain}

%\def\section#1{\@@section{#1}\hskip 5cm}
\begin{document}
	\textbf{\Large{\qquad \qquad \qquad \quad \quad 郭子剑}}
	\newline

	\begin{resume}
		\section{\sc 联系方式}
		\vspace{.05in}

        \begin{tabular}{@{}p{2.5in}p{4in}}
			数据科学研究中心    & zijguo@zju.edu.cn                         \\
			浙江大学，中国 & \url{https://statweb.rutgers.edu/zijguo/} \\
		\end{tabular}

		\section{\sc 教育背景}
		{统计学博士}, 宾夕法尼亚大学沃顿商学院 \hfill 2017
		\\ {\it 博士生导师:} T. Tony Cai\\
		{\it ``Statistical Inference For High-Dimensional Linear Models''}
		
		\vspace*{-3mm}

		{数学理学学士}, 香港中文大学 \hfill 2012\\
		{\it 一级荣誉学士学位}


		\vspace*{-0mm}

        \section{\sc 任职经历}
		{\it 求是讲席教授} \hfill 2025年8月 - 至今 \\ {数据科学研究中心}\\
		浙江大学

		{\it 副教授} (终身教职) \hfill 2022年7月 - 2025年7月 \\
		{统计系}\\ 罗格斯大学——新泽西州立大学

		{\it 助理教授} \hfill 2017年9月 - 2022年6月 \\
		{统计系}\\ 罗格斯大学——新泽西州立大学

		\vspace*{-0mm}

		\section{\sc 研究领域}
		优化与统计、多源学习、非标准推断、因果推断、高维推断，以及在遗传学和健康研究中的应用。

		\section{\sc 荣誉奖励}

        $\bullet$ ICSA 杰出青年学者奖 (Outstanding Young Researcher Award) \hfill{2023年7月}

        \vspace{-3mm}

        $\bullet$ 伯努利学会 (Bernoulli Society) 新锐学者奖荣誉提名 \hfill{2022年11月}

        \vspace{-3mm}

        $\bullet$ ICSA 新锐学者奖 (New Researcher Award) \hfill{2019年12月}

        \vspace{-3mm}

        $\bullet$ IMS 差旅奖 (Travel Award), JSM \hfill{2017年8月}

        \vspace{-3mm}

        $\bullet$ 校长 Gutmann 领导力奖 (President Gutmann Leadership Award), 宾夕法尼亚大学 \hfill{2017年4月}

        \vspace{-3mm}

        $\bullet$ J. Parker Bursk 奖 \hfill{2016年9月}
        \begin{itemize}
            \item[] {\it 由沃顿商学院颁发，以表彰在研究方面的卓越表现。}
        \end{itemize}
        \vspace{-3mm}

        $\bullet$ 流行病学统计青年学者奖 (Statistics in Epidemiology Young Investigator Award), JSM \hfill{2013年8月}
        \begin{itemize}
            \item[] {\it 由美国统计学会 (ASA) 流行病学统计分会颁发。}
        \end{itemize}

        \section{\sc 学术访问}

        {\it 特邀研究员 } (2018年11月, 2023年9-10月, 2024年4-5月, 2025年6月)\\ 
        苏黎世联邦理工学院数学研究所  (邀请人: Peter B{\"u}hlmann)\\
        \vspace*{-6mm}

        {\it 特邀研究员 } (2024年4-5月, 2024年12月, 2025年6月)\\ 
        法国国家信息与自动化研究所, 巴黎  (邀请人: Francis Bach)\\
        \vspace*{-6mm}

        {\it 特邀研究员} (2024年6月)\\ 
        牛津大学统计/经济系  (邀请人: Frank Windmeijer)\\
        \vspace*{-6mm}

        {\it 特邀研究员} (2024年4月)\\ 
        剑桥大学数学系  (邀请人: Richard Samworth)\\
        \vspace*{-6mm}

        {\it 特邀研究员 } (2024年3月)\\ 
        特文特大学(邀请人: Johannes Schmidt-Hieber)\\
        \vspace*{-6mm}

        {\it 特邀研究员 } (2019年9月, 2023年9月, 2024年2-3月)\\ 
        哈佛大学陈曾熙公共卫生学院 (邀请人: Tianxi Cai)\\
        \vspace*{-6mm}

        {\it 特邀研究员} (2024年2月)\\ 
        耶鲁大学经济系  (邀请人: Xiaohong Chen)\\
        \vspace*{-6mm}

        {\it 特邀研究员 } (2024年1-2月)\\ 
        卡内基梅隆大学统计与数据科学系 (邀请人: Larry Wasserman)\\
        \vspace*{-6mm}

        {\it 特邀研究员} (2017年8月) \\ 
        宾夕法尼亚大学佩雷尔曼医学院  (邀请人: Hongzhe Li)\\

		\vspace*{-0mm}

		\section{\sc 代表性论文}
		* 表示作者按姓氏首字母排序；\underline{下划线}表示指导学生；$\#$ 表示同等贡献；\Envelope 表示共同通讯作者。

		\begin{enumerate}
			\setlength{\itemsep}{3pt}

			\item[(R1)] {*}Cai, T. T., and {\bf Guo, Z.} (2017). Confidence intervals for
				high-dimensional linear regression: Minimax rates and adaptivity. {\it Annals of Statistics},
				45(2), 615-646.

			\item[(R2)] {\bf Guo, Z.}, Kang, H., Cai, T. T., and Small, D. S. (2018). Confidence
				Interval for Causal Effects with Invalid Instruments using Two-Stage Hard
				Thresholding.
				{\it Journal of the Royal Statistical Society: Series B (Statistical Methodology)},
				80(4), 793-815.

			\item[(R3)] {\bf Guo, Z.}$^{\#}$, \'{C}evid, D.$^{\#}$, and B\"{u}hlmann, P.
				(2022). Doubly Debiased Lasso: High-Dimensional Inference under Hidden
				Confounding. {\it Annals of Statistics}, 50(3), 1320-1347.

			\item[(R4)] {\bf Guo, Z.} (2023). Causal Inference with Invalid
				Instruments: Post-selection Problems and A Solution Using Searching and Sampling.
				{\it Journal of the Royal Statistical Society: Series B (Statistical Methodology)},
				85(3), 959-985.

			\item[(R5)] {\bf Guo, Z.} (2024). Statistical Inference for Maximin
				Effects: Identifying Stable Associations across Multiple Studies.
				{\it Journal of the American Statistical Association}, 119(547), 1968-1984.

			\item[(R6)] \underline{Yao, M.}, Miller, G., Vardarajan, B., Baccarelli, A.,
				{\bf Guo, Z.} \Envelope, and Liu, Z. \Envelope. (2024+). Robust
				Mendelian Randomization Analysis by Automatically Selecting Valid
				Genetic Instruments with Applications to Identify Plasma Protein
				Biomarkers for Alzheimer{\textquoteright}s Disease. {\it Cell Genomics},
				2024 Dec 11;4(12):100700.

			\item[(R7)] \underline{Wang, Z.}$^{\#}$, Hu, Y.$^{\#}$, Bühlmann, P.\Envelope,
				\& {\bf Guo, Z.}\Envelope (2024). Causal Invariance Learning via
				Efficient Optimization of a Nonconvex Objective. {\it arXiv preprint arXiv:2412.11850.}

			\item[(R8)] {\bf Guo, Z.}, \underline{Wang, Z.}, Hu, Y., and Bach, F. (2025).
				Statistical Inference for Conditional Group Distributionally Robust
				Optimization With Cross-Entropy Loss. {\it arXiv preprint arXiv:2507.09905}

			\item[(R9)] \underline{Zheng, M.}, Bonvini, M.\Envelope, and {\bf Guo, Z.}\Envelope
				(2025). Perturbed Double Machine Learning: Nonstandard Inference Beyond
				the Parametric Length. {\it arXiv preprint arXiv:2511.01222}
		\end{enumerate}

		\section{\sc 预印本}
		
		* 表示作者按姓氏首字母排序；\underline{下划线}表示指导学生；$\#$ 表示同等贡献；\Envelope 表示共同通讯作者。
		\vspace*{2mm}
		\begin{enumerate}
			\setlength{\itemsep}{3pt}

			\item Heng, S.\Envelope, Shen, Y., Guo, Z.\Envelope (2026). Propensity
				Score Propagation: A General Framework for Design-Based Inference with Unknown
				Propensity Scores. {\it arXiv preprint arXiv:2601.13150}

			\item Xiong, X., Guo, Z., Zhu, H., Hong, C., Smoller, J. W., Cai, T., Liu,
				M. (2026). Adversarial Drift-Aware Predictive Transfer: Toward Durable
				Clinical AI. {\it arXiv preprint arXiv:2601.11860}

			\item \underline{Koo, T.} and {\bf Guo, Z.} (2025). Distributionally Robust
				Synthetic Control: Ensuring Robustness Against Highly Correlated
				Controls and Weight Shifts. {\it arXiv preprint arXiv:2511.02632}

			\item \underline{Zheng, M.}, Bonvini, M.\Envelope, and {\bf Guo, Z.}\Envelope
				(2025). Perturbed Double Machine Learning: Nonstandard Inference Beyond
				the Parametric Length. {\it arXiv preprint arXiv:2511.01222}

			\item {\bf Guo, Z.}, \underline{Wang, Z.}, Hu, Y., and Bach, F. (2025).
				Statistical Inference for Conditional Group Distributionally Robust
				Optimization With Cross-Entropy Loss. {\it arXiv preprint arXiv:2507.09905}

			\item \underline{Wang, Z.}, Liu, M., Lei, J., Bach, F.\Envelope, \& {\bf Guo, Z.}\Envelope
				(2025). StablePCA: Learning Shared Representations across Multiple
				Sources via Minimax Optimization. {\it arXiv preprint arXiv:2505.00940.}

			\item \underline{Scheidegger, C.}, {\bf Guo, Z.}, \& Bühlmann, P. (2025). Inference
				for Heterogeneous Treatment Effects with Efficient Instruments and Machine
				Learning. {\it arXiv preprint\\ arXiv:2503.03530.}

			\item \underline{Gu, Y.}, Fang, C., Xu, Y., {\bf Guo, Z.}, \& Fan, J. (2025).
				Fundamental Computational Limits in Pursuing Invariant Causal Prediction
				and Invariance-Guided Regularization.
				{\it arXiv preprint arXiv:2501.17354.}

			\item \underline{Wang, Z.}$^{\#}$, Hu, Y.$^{\#}$, Bühlmann, P., \& {\bf Guo, Z.}
				(2024). Causal Invariance Learning via Efficient Optimization of a
				Nonconvex Objective. {\it arXiv preprint arXiv:2412.11850.}

			\item \underline{Rakshit, P.}, \& {\bf Guo, Z.} (2024). Statistical
				Inference in High-dimensional Poisson Regression with Applications to
				Mediation Analysis. {\it arXiv preprint arXiv:2410.20671.}

			\item \underline{Wang, Z.}, Si, N., {\bf Guo, Z.}, and Liu, M. (2024).
				Multi-source Stable Variable Importance Measure via Adversarial Machine
				Learning. {\it arXiv preprint arXiv:2409.07380.}

			\item \underline{Zhan, K.}, \underline{Xiong, X.}, {\bf Guo, Z.}, Cai, Tianxi,
				and Liu, M. (2024). Transfer Learning Targeting Mixed Population: A
				Distributional Robust Perspective. {\it arXiv preprint arXiv:2407.20073.}

			\item {*}Fan, Q., {\bf Guo, Z.}, \underline{Mei, Z.}, and Zhang, C. (2023).
				Uniform Inference for Nonlinear Endogenous Treatment Effects with High-Dimensional
				Covariates. {\it arXiv preprint arXiv:2310.08063}.

			\item \underline{Xiong, X.}, {\bf Guo, Z.}\Envelope, and Cai, Tianxi
				\Envelope(2023). Distributionally Robust Transfer Learning. {\it arXiv preprint arXiv:2309.06534}.

			\item \underline{Liu, Y.}, Liu, M., {\bf Guo, Z.}, and Cai, Tianxi. (2023).
				Surrogate-Assisted Federated Learning of high dimensional Electronic
				Health Record Data. {\it arXiv preprint arXiv:2302.\\04970}.

			\item {\bf Guo, Z.}, \underline{Zheng, M.}, and B\"{u}hlmann, P. (2022).
				Robustness Against Weak or Invalid Instruments: Exploring Nonlinear Treatment
				Models with Machine Learning. {\it arXiv preprint arXiv:2203.12808}.

			\item {*}{\bf Guo, Z.}, \underline{Yuan, W.} and Zhang, C. (2019). Decorrelated
				Local Linear Estimator: Inference for Non-linear Effects in High-dimensional
				Additive Models. {\it arXiv preprint arXiv:1907.12732}.
				{Reject and Resubmit} at {\it Journal of Machine Learning Research}.
		\end{enumerate}
		\section{\sc 出版物}
		* 表示作者按姓氏首字母排序；\underline{下划线}表示指导学生；$\#$ 表示同等贡献；\Envelope 表示共同通讯作者。
		\vspace{2mm}
		\begin{enumerate}
			\setcounter{enumi}{0}
			\setlength{\itemsep}{3pt}

			\item \underline{Wang, Z.}, Bühlmann, P., and {\bf Guo, Z.}(2025+). Distributionally
				Robust Learning for Multi-source Unsupervised Domain Adaptation.
				{\it Annals of Statistics}, to appear.

			\item {\bf Guo, Z.}, Li, X., Han, L., and Cai, Tianxi. (2025+). Robust
				Inference for Federated Meta-Learning. {\it Journal of the American Statistical Association},
				to appear.

			\item \underline{Chang, T. H.}, {\bf Guo, Z.}, and Malinsky, D. (2025+).
				Post-selection inference for causal effects after causal discovery. {\it Biometrika},
				to appear.

			\item \underline{Scheidegger, C.}, {\bf Guo, Z.}, and B\"{u}hlmann, P. (2024+).
				Spectral Deconfounding for High-Dimensional Sparse Additive Models.
				{\it ACM/IMS Journal of Data Science}, to appear.

			\item \underline{Yao, M.}, Miller, G., Vardarajan, B., Baccarelli, A.,
				{\bf Guo, Z.} \Envelope, and Liu, Z \Envelope. (2024+). Robust Mendelian
				Randomization Analysis by Automatically Selecting Valid Genetic Instruments
				with Applications to Identify Plasma Protein Biomarkers for Alzheimer{\textquoteright}s
				Disease. {\it Cell Genomics}, to appear.

			\item {Lin, Y.}, {\bf Guo, Z.}, Sun, B., and Lin, Z. (2024+). Testing High-Dimensional
				Mediation Effect with Arbitrary Exposure-Mediator Coefficients. {\it Test},
				to appear.

			\item Carl, D., Emmenegger, C., Bühlmann, B., {\bf Guo, Z.} (2024+). TSCI:
				Two Stage Curvature Identification for Causal Inference with Invalid Instruments.
				{\it Journal of Statistical Software}, to appear.

			\item {*}Fan, Q., {\bf Guo, Z.}, and \underline{Mei, Z.} (2024+). A Heteroskedasticity-Robust
				Overidentifying Restriction Test with High-Dimensional Covariates.
				{\it Journal of Business \& Economic Statistics}, to appear.

			\item {\bf Guo, Z.} (2024). Statistical Inference for Maximin Effects: Identifying
				Stable Associations across Multiple Studies.
				{\it Journal of the American Statistical Association}, 119(547), 1968-1984.

			\item Kang, H., {\bf Guo, Z.}, Liu, Z., and Small, D. (2024). Identification
				and Inference with Invalid Instruments.
				{\it Annual Review of Statistics and Its Application}, to appear.

			\item Ma, R., {\bf Guo, Z.}, Cai, T. T., and Li, H. (2024). Statistical
				Inference of Genetic Relatedness using High-Dimensional Logistic
				Regression. {\it Statistica Sinica}, 34 (2024): 1023-1043.. %{\it arXiv preprint arXiv:2202.10007}.

			\item \underline{Rakshit, P.}, \underline{Wang, Z.}, Cai, T. T., and {\bf Guo, Z.}
				(2024). \texttt{SIHR}: An R Package for Statistical Inference in High-dimensional
				Linear and Logistic Regression Models. {\it R Journal}, to appear.

			\item *Cai, T. T., {\bf Guo, Z.}, and Xia, Y. (2023). Statistical
				inference and large-scale multiple testing for high-dimensional
				regression models. {\it Test (with discussion)}, 32(4), 1135-1171.

			\item {\bf Guo, Z.} (2023). Causal Inference with Invalid Instruments: Post-selection
				Problems and A Solution Using Searching and Sampling.
				{\it Journal of the Royal Statistical Society: Series B (Statistical Methodology)},
				85(3), 959-985.

			\item {*}Cai, T. T., {\bf Guo, Z.}, and Ma, R. (2023). Statistical
				Inference for High-Dimensional Generalized Linear Models with Binary
				Outcomes. {\it Journal of the American Statistical Association}, 118 (542),
				1319-1332.

			\item \underline{Koo, T.}, {Lee, Y.}, Small, D. S., and {\bf Guo, Z.} (2023).
				\texttt{RobustIV} and \texttt{controlfunctionIV}: Robustiv and
				controlfunctioniv: Causal inference for linear and nonlinear models with
				invalid instrumental variables. {\it Observational Studies}, 9(4), 97-120.

			\item Hou, J., {\bf Guo, Z.}, and Cai, T. (2023). Surrogate assisted semi-supervised
				inference for high dimensional risk prediction. {\it Journal of Machine Learning Research},
				24(265), 1-58.
			
			\item Wang, X., Zhou, H., $\cdots$, 4CE, Avillach, P.\Envelope, {\bf Guo, Z.}\Envelope,
				and Cai, Tianxi\Envelope. (2022) SurvMaximin: Robust Federated Approach to
				Transporting Survival Risk Prediction Models.
				{\it Journal of Biomedical Informatics}, 134 (2022): 104-176. 

			\item {\bf Guo, Z.}, \'{C}evid, D., and B\"{u}hlmann, P. (2022). Doubly Debiased
				Lasso: High-Dimensional Inference under Hidden Confounding. {\it Annals of Statistics},
				50 (3), 1320 - 1347. 

			\item {*}{\bf Guo, Z.} and Zhang, C. (2022). Extreme Nonlinear Correlation
				for Multiple Random Variables and Stochastic Processes with Applications
				to Additive Models. {\it Stochastic Processes and their Applications}. 150,
				1037-1058. 

			\item {\bf Guo, Z.}, Renaux, C., B\"{u}hlmann, P., and Cai, T. T. (2021). Group
				Inference in High Dimensions with Applications to Hierarchical Testing.
				{\it Electronic Journal of Statistics}, 15(2), 6633-6676. %{arXiv preprint arXiv:1909.01503}.

			\item {\bf Guo, Z.}, \underline{Rakshit, P.}, Herman, D., and Chen, J. (2021).
				Inference for Case Probability in High-dimensional Logistic Regression. {\it Journal of Machine Learning Research},
				22(254), 1-54. %{arXiv preprint arXiv:2012.07133}.

			\item {*}Cai, Tianxi, Cai, T. T., and {\bf Guo, Z.} (2021). Optimal Statistical
				Inference for Individualized Treatment Effects in High-dimensional Models.
				{\it Journal of the Royal Statistical Society: Series B (Statistical Methodology)},
				83(4), 669-719. %{arXiv preprint arXiv:1904.12891}.

			\item {*}Cai, T. T., and {\bf Guo, Z.} (2020). Semi-supervised Inference for
				Explained Variance in High-dimensional Linear Regression and Its Applications.
				{\it Journal of the Royal Statistical Society: Series B (Statistical Methodology)},
				82(2), 391-419. %{arXiv preprint arXiv:1806.06179}.

			\item {\bf Guo, Z.}, Wang, W., Cai, T. T., and Li, H. (2019). Optimal
				Estimation of Genetic Relatedness in High-dimensional Linear Models. {\it Journal of the American Statistical Association},
				114(525), 358-369. %{arXiv preprint arXiv:1605.07244}.

			\item {\bf Guo, Z.}, Kang, H., Cai, T. T., and Small, D. S. (2018). Testing
				Endogeneity with High Dimensional Covariates.
				{\it The Journal of Econometrics}, 207(1), 175-187. %{arXiv preprint arXiv:1609.06713}.

			\item {\bf Guo, Z.}, Kang, H., Cai, T. T., and Small, D. S. (2018). Confidence
				Interval for Causal Effects with Invalid Instruments using Two-Stage Hard
				Thresholding.
				{\it Journal of the Royal Statistical Society: Series B (Statistical Methodology)},
				80(4), 793-815. %{arXiv preprint arXiv:1603.05224}.% representative

			\item {*}Cai, T. T., and {\bf Guo, Z.} (2018). Accuracy assessment for high-dimensional
				linear regression. {\it Annals of Statistics}, 46(4), 1807-1836. %{arXiv preprint arXiv:1603.03474}. % representative

			\item {\bf Guo, Z.}, Small, D. S., Gansky, S. A., and Cheng, J. (2018). Mediation
				analysis for count and zero-inflated count data without sequential ignorability
				and its application in dental studies.
				{\it Journal of the Royal Statistical Society: Series C (Applied Statistics)},
				67(2), 371-394.

			\item Cheng, J., Cheng N. F., {\bf Guo, Z.}, Gregorich, S., Amid I. I.,
				and Gansky, S. A. (2018). Mediation analysis for count and zero-inflated
				count data. {\it Statistical Methods in Medical Research}, 27(9), 2756-2774. %{arXiv preprint arXiv:1601.06294}.

			\item {*}Cai, T. T., and {\bf Guo, Z.} (2017). Confidence intervals for high-dimensional
				linear regression: Minimax rates and adaptivity. {\it Annals of Statistics},
				45(2), 615-646. %{arXiv preprint arXiv:1506.05539}. % representative

			\item {\bf Guo, Z.}, and Small, D. S. (2016). Control function
				instrumental variable estimation of nonlinear causal effect models. {\it Journal of Machine Learning Research},
				17(100), 1-35. %{arXiv preprint arXiv:1602.01051}.

			\item {\bf Guo, Z.}, Cheng, J., Lorch, S. A., and Small, D. S. (2014). Using
				an instrumental variable to test for unmeasured confounding.
				{\it Statistics in Medicine}, 33(20), 3528-3546. %{arXiv preprint arXiv:1601.06288}.

			\item {\bf Guo, Z.}, Kogan, R., Qiu, H., and Strichartz, R. S. (2014).
				Boundary value problems for a family of domains in the Sierpinski gasket.
				{\it Illinois Journal of Mathematics}, 58(2), 497-519. %{arXiv preprint arXiv:1310.6463}.
		\end{enumerate}

		\section{\sc 软件}
		\begin{enumerate}
            \setlength{\itemsep}{3pt}

            \item {R 包} \texttt{\bf SIHR}: statistical inference in high-dimensional
                regression. \\ 地址: \url{https://cran.r-project.org/web/packages/SIHR/index.html}

            \item {R 包} \texttt{\bf RobustIV}: robust causal inference with possibly
                invalid instruments. \\ 地址: \url{https://cran.r-project.org/web/packages/RobustIV/index.html}

            \item {R 包} \texttt{\bf TSCI}: two stage curvature identification with
                machine learning. \\ 地址: \url{https://cran.r-project.org/web/packages/TSCI/index.html}

            \item {R 包} \texttt{\bf maczic}: mediation analysis for count and zero-inflated
                count data. \\ 地址: \url{https://cran.r-project.org/web/packages/maczic/index.html}

            \item {R 包} \texttt{\bf MaximinInfer}: inference for maximin effects
                in high-dimensional settings. \\ 地址: \url{https://cran.r-project.org/web/packages/MaximinInfer/index.html}

            \item {R 包} \texttt{\bf DLL}: inference for function derivative in high-dimensional
                additive models. \\ 地址: \url{https://cran.r-project.org/web/packages/DLL/index.html}

            \item {R 包} \texttt{\bf DDL}: inference for regression parameters in
                high-dimensional models with hidden confounders. \\ 地址: \url{https://cran.r-project.org/web/packages/DDL/index.html}
            
            \item {R 包} \texttt{\bf controlfunctionIV}: control function method with
                possibly invalid instruments. \\ 地址: \url{https://cran.r-project.org/web/packages/controlfunctionIV/index.html}
            
        \end{enumerate}
			
		R 代码均可在 \url{https://github.com/zijguo} 找到.

		\section{\sc 过往科研项目}
        \begin{enumerate}
            \setlength{\itemsep}{3pt}

            \item 美国国立卫生研究院\\ R01AG086379 ``Robust Mendelian
                Randomization Framework with Multi-Omics Data for Alzheimer's Disease
                and Related Dementias"
                \begin{itemize}
                    \item[-] {职责: 共同项目负责人 (Co-PI)} (负责人: Dr. Zhonghua Liu)
                \end{itemize}

            \item 美国国立卫生研究院\\ R01LM013614 ``Semi-supervised
                Approaches to Denoising Electronic Health Records Data for Risk
                Prediction"
                \begin{itemize}
                    \item[-] {职责: 共同项目负责人 (Co-PI)} (负责人: Dr. Tianxi Cai)
                \end{itemize}

            \item 美国国立卫生研究院\\ R01GM140463 ``Predictive Modeling
                with High-Dimensional Incomplete Data."
                \begin{itemize}
                    \item[-] {职责: 项目负责人 (PI)} (共同负责人: Dr. Jinbo Chen)
                \end{itemize}

            \item 美国国家科学基金会\\ DMS 2015373 ``Repro Sampling Method: A
                Transformative Artificial-Sample-Based Inferential Framework with Applications
                to Discrete Parameter, High-Dimensional Data, and Rare Events Inferences."
                \begin{itemize}
                    \item[-] {职责: 共同项目负责人 (Co-PI)} (负责人: Dr. Min-ge Xie)
                \end{itemize}

            \item 美国国家科学基金会\\ DMS 1811857 ``Inference in High-Dimensional
                Linear Models: Methods, Theory and Applications."
                \begin{itemize}
                    \item[-] {职责: 项目负责人 (PI)}
                \end{itemize}

            \item 美国国立卫生研究院\\ R56-HL-138306-01 ``Statistics Methods
                for Analyzing Electronic Health Record Data."
                \begin{itemize}
                    \item[-] {职责: 共同研究员 (Co-Investigator)} (负责人: Dr. Jinbo Chen)
                \end{itemize}

            \item 宾夕法尼亚大学医学院\\ ``Statistics Methods for
                Analyzing Electronic Health Record Data."
                \begin{itemize}
                    \item[-] {职责: 高级研究员 (Senior Investigator)} (负责人: Dr. Jinbo Chen)
                \end{itemize}
        \end{enumerate}

		\section{\sc 授课经历}
		{\bf 授课讲师}

        \vspace*{1mm}

        \begin{itemize}
            \item[-] {\bf 罗格斯大学(博士生课程)}

            \item[] STAT 593: Theory of Statistics \hfill{2024年秋季}

            \item[] STAT 593: Theory of Statistics \hfill{2023年秋季}

            \item[] STAT 594: Advanced Modern Statistical Inference II \hfill{2019年春季}

            \item[] {\it 授课评分 (Instructor Rating):} 4.82 / 5.0

            \item[-] {\bf 罗格斯大学(研究生课程)}

            \item[] FSRM 588: Financial Data Mining \hfill{2025年春季}

            \item[] FSRM 588: Financial Data Mining \hfill{2021年秋季}

            \item[] {\it 授课评分 (Instructor Rating):} 4.75 / 5.0

            \item[] FSRM 588: Financial Data Mining (线上课程) \hfill{2021年春季}

            \item[] FSRM 588: Financial Data Mining \hfill{2020年春季}

            \item[] {\it 授课评分 (Instructor Rating):} 4.50 / 5.0

            \item[] FSRM 588: Financial Data Mining \hfill{2019年秋季}

            \item[] {\it 授课评分 (Instructor Rating):} 4.75 / 5.0

            \item[] FSRM 588: Financial Data Mining \hfill{2018年秋季}

            \item[] {\it 授课评分 (Instructor Rating):} 4.71 / 5.0

            \item[] FSRM 588: Financial Data Mining \hfill{2017年秋季}

            \item[] {\it 授课评分 (Instructor Rating):} 4.82 / 5.0

            \item[-] {\bf 罗格斯大学(本科生课程)}

            \item[] STAT 384: Intermediate Statistical Analysis (线上课程) \hfill{2021年春季}

            \item[] STAT 484: Basic Applied Statistics (线上课程) 
        \end{itemize}
        \vspace*{1mm}
        \begin{itemize}
            \item[-] {\bf 宾夕法尼亚大学沃顿商学院}

            \item[] STAT 111 : Introductory Statistics \hfill{2016年夏季}
            
        \end{itemize}

        \vspace{-3mm}

        {\bf 习题课讲师} \hfill{2014年秋季}
        \begin{itemize}
            \item[] {\bf 宾夕法尼亚大学沃顿商学院}

            \item[] STAT 111: Introductory Statistics
        \end{itemize}

        \vspace{-3mm}

        {\bf 助教}
        \begin{itemize}
            \item[] {\bf 宾夕法尼亚大学沃顿商学院} 

            \item[] STAT 102: Business Statistics \hfill{2017年春季}

            \item[] STAT 970: Mathematical Statistics \hfill{2016年秋季}

            \item[] STAT 622: Statistical Modeling \hfill{2016年春季}

            \item[] STAT 550: Mathematical Statistics \hfill{2015年秋季}
        \end{itemize}

		\vspace*{-3mm}

		\section{\sc 学生指导}

        {\bf 博士生导师:}
        \vspace{2mm}
        \begin{enumerate}
            \item Prabrisha Rakshit (2023年毕业；Assistant Professor in the Department
                of Operations Management, Quantitative Methods, and Information Systems at
                the Indian Institute of Management Udaipur);

            \item Taehyeon Koo (2025年毕业；与 Dr. Nicole Pashley 联合指导；
                首职: Postdoctoral fellow at Columbia University);

            \item Zhenyu Wang (罗格斯大学, 预计2027年毕业);

            \item Mengchu Zheng (罗格斯大学, 预计2028年毕业);

            \item Chao Qin (浙江大学)

            \item Yuebei Shi (浙江大学)

            \item Chen Ling (浙江大学)

            \item Jiahang Shao (浙江大学)

            \item Wenyuan Shu (浙江大学)
        \end{enumerate}

        {\bf 博士后合作导师:}
        \vspace{2mm}
        \begin{enumerate}
            \item Yichuan Bai (Ph.D. at Iowa State)
        \end{enumerate}

		\section{\sc 学术服务}

        \begin{enumerate}
            \item 副主编 (Associate Editor)
                \begin{itemize}
                    \item {\it Journal of the American Statistical Association}, 2023- 至今
                \end{itemize}

            \item 因果推断阅读小组组织者 (Organizer of Causal Inference Reading Group) \\ 
                (与 Nicole Pashley 和 Tirthankar Dasgupta 共同组织), 罗格斯大学统计系

            \item 系年度研讨会主席 (Department Retreat Chair, 2019-2020), 罗格斯大学统计系

            \item 系学术报告会主席 (Department Seminar Chair, 2018-2019), 罗格斯大学统计系

            \item 其他罗格斯大学委员会服务:
                \begin{itemize}
                    \item 系网站委员会 (Department website committee, 2020-2023)

                    \item 系年度研讨会委员会 (Department retreat committee, 2017-2018)

                    \item FSRM/MSDS 项目委员会 (FSRM/MSDS committee, 2017-2023)

                    \item 博士生考试委员会 (Ph.D. Exam committee, 2018-2020)

                    \item 研究生课程委员会 (Graduate Curriculum committee, 2019-2021)
                \end{itemize}

            \item 2019年罗格斯统计研讨会组织委员会 (Organizing Committee for 2019 Rutgers Statistics Symposium)

            \item ICSA 2019 第11届国际会议程序委员会 (Program Committee for ICSA 2019 11th International Conference)

            \item 2018年 ICSA 应用研讨会地方组织委员会 (Local Organizing Committee for 2018 ICSA Applied Symposium)

            \item 期刊审稿人 (Reviewer for the following journals): {\it Annals of Statistics}, {\it Journal of the American Statistical Association},
                {\it Journal of the Royal Statistical Society: Series B (Statistical Methodology)},
                {\it Biometrika}, {\it Journal of Machine Learning}, {\it Journal of Econometrics},
                {\it Biometrics},{\it Statistica Sinica}, {\it IEEE International Symposium on Information Theory},
                {\it Journal of Applied Statistics}, {\it COLT}.
        \end{enumerate}

		\vspace*{1mm}

		\section{\sc 短期课程}

		\begin{enumerate}
            \setlength{\itemsep}{3pt}

            \item {\it ``Nonregular Inference"}, 耶鲁大学经济系, 
                2024年2月

            \item {\it ``Nonregular Inference"}, CELEHS Distinguished Lecture Series, 
            哈佛大学生物统计系, 2024年2月

            \item {\it ``Robust Deconfounding"}, 哈佛大学生物统计系, 
                2024年2月
        \end{enumerate}

		\section{\sc 学术报告}

		\begin{enumerate}
			\setlength{\itemsep}{3pt}

			\item 系学术报告, 牛津大学经济系, 英国, {\it ``Adversarially Robust Learning: Identification, Estimation, and Uncertainty Quantification.''}, 
                2024年6月

            \item 系学术报告, 牛津大学统计系, 英国, {\it ``Robust Instrumental Variable Analysis.''}, 
                2024年6月

            \item 系学术报告, 日内瓦大学, 瑞士, {\it ``Adversarially Robust Learning: Identification, Estimation, and Uncertainty Quantification.''}, 
                2024年5月

            \item 系学术报告, 洛桑联邦理工学院, 瑞士, {\it ``Adversarially Robust Learning: Identification, Estimation, and Uncertainty Quantification.''}, 
                2024年5月

            \item 系学术报告, 苏黎世联邦理工学院统计研讨会, 瑞士, {\it ``Adversarially Robust Learning: Identification, Estimation, and Uncertainty Quantification.''}, 
                2024年5月

            \item 系学术报告, 法国国家信息与自动化研究所, 法国, {\it ``Adversarially Robust Learning: Identification, Estimation, and Uncertainty Quantification.''}, 
                2024年4月

            \item 学术报告, 阿姆斯特丹机器学习实验室, {\it ``Adversarially Robust Learning: Identification, Estimation, and Uncertainty Quantification.''}, 
                2024年3月

            \item 系学术报告, 耶鲁大学经济系, 美国, {\it ``Robust Instrumental Variable Analysis.''}, 
                2024年2月

            \item 因果推断阅读小组报告, 卡内基梅隆大学统计与数据科学系, 美国, {\it ``Robust Deconfounding.''}, 
                2024年1月

            \item 机器学习阅读小组报告, 卡内基梅隆大学统计与数据科学系, 美国, {\it ``Adversarially Robust Learning: Identification, Estimation, and Uncertainty Quantification.''}, 
                2024年1月

            \item 系学术报告, 香港理工大学统计系, 香港, {\it ``Domain Generalization with Adversarially Robust Learning: Identification, Estimation, and Uncertainty Quantification.''}, 
                2024年1月

            \item 系学术报告, 香港中文大学统计系, 香港, {\it ``Domain Generalization with Adversarially Robust Learning: Identification, Estimation, and Uncertainty Quantification.''}, 
                2024年1月

			\item 特邀报告, ICSDS, 葡萄牙里斯本, {\it ``Distributionally Robust Machine Learning with Multi-source Data.''}, 
                2023年9月

            \item 系学术报告, 新泽西理工学院数学系, 美国, {\it ``Two Stage Curvature Identification with Machine Learning: Causal Inference with Possibly Invalid Instrumental Variables.''}, 
                2023年10月

            \item ETH/UZH 统计与数据科学研究系列研讨会, 瑞士苏黎世, 
                {\it ``Robust Causal Inference with Possibly Invalid Instruments: Post-selection Problems and A Solution Using Searching and Sampling.''}, 
                2023年9月

            \item 哈佛大学 Cai 实验室研讨会, 哈佛大学陈曾熙公共卫生学院, 美国, {\it ``Distributionally Robust Machine Learning with Multi-source Data.''}, 
                2023年9月

            \item 系学术报告, 香港科技大学资讯、商业统计及运营管理系, 中国香港, 
                {\it ``Robust Inference for Federated Meta-Learning.''}, 2023年8月

            \item 特邀报告, 2023 杭州数据科学会议, 中国杭州, 
                {\it ``Statistical Inference for Maximin Effects: Identifying Stable Associations across Multiple Studies.''}, 
                2023年8月

            \item 特邀演讲 (线上), EcoSta 2023, 日本东京, {\it ``Two Stage Curvature Identification with Machine Learning: Causal Inference with Possibly Invalid Instrumental Variables.''}, 
                2023年8月

            \item 特邀演讲, 第九届人大国际统计论坛 (9th RUC-IFS), 中国北京, {\it ``Statistical Inference for Maximin Effects: Identifying Stable Associations across Multiple Studies.''}, 
                2023年7月

            \item 特邀演讲, 联合数据科学统计会议 (JCSDS), 中国北京, {\it ``Robust Inference for Federated Meta-Learning.''}, 
                2023年7月

            \item 特邀演讲, ICSA 2023, 中国香港, {\it ``Causal Inference with Invalid Instruments: Post-selection Problems and A Solution Using Searching and Sampling.''}, 
                2023年7月

            \item 特邀报告 (线上), ENSAE/巴黎综合理工学院联合统计研讨会, 法国, 
                {\it ``Doubly Debiased Lasso: High-Dimensional Inference under Hidden Confounding.''}, 
                2023年6月

            \item 特邀报告 (线上), 法国国家信息与自动化研究所 (Inria), 法国, 
                {\it ``Statistical Inference for Maximin Effects: Identifying Stable Associations across Multiple Studies.''}, 
                2023年5月

            \item 特邀报告 (线上), 北京师范大学, 中国, 
                {\it ``Statistical Inference for Maximin Effects: Identifying Stable Associations across Multiple Studies.''}, 
                2023年5月

            \item 特邀报告 (线上), 北京大学, 中国, 
                {\it ``Robust Inference for Federated Meta-Learning.''}, 2023年3月

            \item 特邀报告 (线上), IMS 新锐学者组、苏黎世青年数据科学研究员研讨会与 YoungStatS 项目联合网络研讨会, 
                {\it ``Statistical Inference for Maximin Effects: Identifying Stable Associations across Multiple Studies.''}, 
                2023年3月

			\item 特邀报告 (线上), CMStatistics 2022, 英国伦敦, {\it ``Doubly Debiased Lasso: High-Dimensional Inference under Hidden Confounding.''}, 
                2022年12月

            \item Levin 讲座, 哥伦比亚大学生物统计系, 美国, 
                {\it ``Multi-data Replicability via the Maximin Criterion: High-Dimensional Inference for Group Distributionally Robust Models.''}, 
                2022年9月

            \item 实验室研讨会 (线上, 由 Tianxi Cai 教授主持), 哈佛大学陈曾熙公共卫生学院, 美国, {\it ``Robust Inference for Federated Meta-Learning.''}, 
                2022年9月

            \item RAND 中心因果推断研讨会 (线上), 
                {\it ``Two Stage Curvature Identification with Machine Learning: Causal Inference with Possibly Invalid Instrumental Variables.''}, 
                2022年8月

            \item 特邀报告人, 2022 联合统计会议 (JSM 2022), 美国, 
                {\it ``Causal Inference with Invalid Instruments: Post-selection Problems and A Solution Using Searching and Sampling.''}, 
                2022年8月

            \item 特邀报告人(线上), ICSA 2022, 佛罗里达大学, 美国, {\it ``Doubly Debiased Lasso: High-Dimensional Inference under Hidden Confounding.''}, 
                2022年6月

            \item 在线因果推断研讨会 (线上, 评论人: Professor Frank Windmeijer), {\it ``Two Stage Curvature Identification with Machine Learning: Causal Inference with Possibly Invalid Instrumental Variables.''}, 
                2022年4月

            \item 因果推断中心研讨会 (线上), 宾夕法尼亚大学, 美国, 
                {\it ``Two Stage Curvature Identification with Machine Learning: Causal Inference with Possibly Invalid Instrumental Variables.''}, 
                2022年3月

            \item 系学术报告 (线上), 威斯康星大学麦迪逊分校生物统计与医学信息学系, 美国, {\it ``Decorrelated Local Linear Estimator: Inference for Non-linear Effects in High-dimensional Additive Models.''}, 
                2022年2月

            \item 实验室研讨会 (Lab seminar, 线上, 由 Tianxi Cai 教授主持), 哈佛大学陈曾熙公共卫生学院, 美国, {\it ``Transfer Learning with Multi-source Data: High-Dimensional Inference for Group Distributionally Robust Models.''}, 
                2022年1月

            \item 特邀报告人 (线上), CMStatistics 2021, {\it ``Inference for Case Probability in High-dimensional Logistic Regression.''}, 
                2021年12月

            \item 特邀报告人 (线上), JSM 2021, 美国西雅图, {\it ``Inference for High-dimensional Maximin Effects in Heterogeneous Regression Models Using a Sampling Approach.''}, 
                2021年8月

            \item 特邀报告人 (线上), 第一届统计及相关领域国际会议, 卢森堡, 
                {\it ``Inference for High-dimensional Maximin Effects in Heterogeneous Regression Models Using a Sampling Approach.''}, 
                2021年7月

            \item 实验室研讨会 (Lab seminar, 线上, 由 Tianxi Cai 教授主持), 哈佛大学 (Harvard University) 陈曾熙公共卫生学院, 美国, {\it ``Maximin Effect and Distributional Robustness: A Review and New Advances.''}, 
                2021年6月

            \item 系学术报告 (线上), 华东师范大学统计系, 中国上海, 
                {\it ``Inference for High-dimensional Maximin Effects in Heterogeneous Regression Models Using a Sampling Approach.''}, 
                2021年5月

            \item 因果推断中心研讨会 (,线上), 宾夕法尼亚大学, 美国, 
                {\it ``Post-selection Problems for Causal Inference with Invalid Instruments: A Solution Using Searching and Sampling.''}, 
                2021年5月

            \item 系学术报告 (线上), 香港中文大学经济系, 中国香港, 
                {\it ``Post-selection Problems for Causal Inference with Invalid Instruments: A Solution Using Searching and Sampling.''}, 
                2021年4月

            \item 系学术报告 (线上), 香港大学统计及精算学系, 中国香港, 
                {\it ``Doubly Debiased Lasso: High-Dimensional Inference under Hidden Confounding.''}, 
                2021年4月

            \item 系学术报告 (线上), 埃克塞特大学医学院, 英国埃克塞特, 
                {\it ``Inference for Non-linear Treatment Effects with Control Function Methods.''}, 
                2021年2月

			\item 特邀报告 (线上), CMStatistics 2020, 英国伦敦, {\it ``Doubly Debiased Lasso: High-Dimensional Inference under Hidden Confounding.''}, 
                2020年12月

            \item 系学术报告 (线上), 康奈尔大学统计系, 美国, {\it ``Doubly Debiased Lasso: High-Dimensional Inference under Hidden Confounding.''}, 
                2020年10月

            \item 特邀报告 (线上), JSM 2020, 美国费城, {\it ``Doubly Debiased Lasso: High-Dimensional Inference under Hidden Confounding and Measurement Errors.''}, 
                2020年8月

            \item 特邀参会者 (线上), {\it ``Mathematical and Statistical Challenges in Uncertainty Quantification.''}, 
                英国剑桥, 2020年7月

            \item 系学术报告 (线上), 加州大学戴维斯分校统计系, 美国, {\it ``Doubly Debiased Lasso: High-Dimensional Inference under Hidden Confounding and Measurement Errors.''}, 
                2020年5月

            \item 因果推断中心研讨会 (线上), 宾夕法尼亚大学, 美国, 
                {\it ``Doubly Debiased Lasso: High-Dimensional Inference under Hidden Confounding and Measurement Errors.''}, 
                2020年5月

            \item 生物统计阅读小组报告 (线上, 由 Jinbo Chen 教授主持), 宾夕法尼亚大学, 美国, {\it ``Group Inference in High Dimensions with Applications to Hierarchical Testing.''}, 
                2020年5月

            \item 系学术报告, 华东师范大学统计系, 中国上海, 
                {\it ``Group Inference in High Dimensions with Applications to Hierarchical Testing.''}, 
                2019年12月

            \item 特邀报告, 第11届 ICSA 国际会议, 中国杭州, {\it ``Group Inference in High Dimensions with Applications to Hierarchical Testing.''}, 
                2019年12月

            \item 特邀报告, 李锡钦教授纪念国际统计会议, 中国香港, {\it ``Group Inference in High Dimensions with Applications to Hierarchical Testing.''}, 
                2019年12月

            \item 因果推断阅读小组报告(由 James Robins 教授主持), 哈佛大学陈曾熙公共卫生学院, 美国, {\it ``Semi-supervised Inference for Explained Variance in High-dimensional Linear Regression and Its Applications.''}, 
                2019年9月

            \item 系学术报告, 华东师范大学统计系, 中国上海, 
                {\it ``Individualized Treatment Selection: A Hypothesis Testing Approach in High-dimensional Models.''}, 
                2019年6月

            \item 特邀报告, 2019 杭州数据科学会议, 中国杭州, 
                {\it ``Local Inference in High-dimensional Sparse Additive Modeling.''}, 
                2019年5月

            \item 系学术报告, 香港城市大学数据科学学院, 中国香港, 
                {\it ``Individualized Treatment Selection: A Hypothesis Testing Approach in High-dimensional Models.''}, 
                2019年5月

            \item 系学术报告, 香港科技大学资讯、商业统计及运营管理系, 中国香港, 
                {\it ``Local Inference in High-dimensional Sparse Additive Modeling.''}, 2019年5月

            \item 系学术报告, 弗吉尼亚大学统计系, 美国, {\it ``Local Inference in High-dimensional Sparse Additive Modeling.''}, 
                2019年3月

            \item 特邀报告, 2019 ICSA 数据科学会议, 中国云南西双版纳, 
                {\it ``Individualized Treatment Selection: A Hypothesis Testing Approach in High-dimensional Models.''}, 
                2019年1月

			\item 青年学者会议 (Young Research Session), Lawrence D. Brown 纪念研讨会, 宾夕法尼亚大学, 美国,
                {\it ``Individualized Treatment Selection: A Hypothesis Testing Approach in High-dimensional Models.''}, 
                2018年11月

            \item 系学术报告, 苏黎世联邦理工学院数学系统计研讨会, 瑞士, 
                {\it ``Semi-supervised Inference for Explained Variance in High-dimensional Linear Regression and Its Applications.''}, 
                2018年11月

            \item 系学术报告, 新泽西理工学院数学系, 美国, {\it ``Semi-supervised Inference for Explained Variance in High-dimensional Linear Regression and Its Applications.''}, 
                2018年11月

            \item 系学术报告, 普林斯顿大学运筹与金融工程系, 美国, {\it ``Semi-supervised Inference for Explained Variance in High-dimensional Linear Regression and Its Applications.''}, 
                2018年10月

            \item 系学术报告, 香港科技大学资讯、商业统计及运营管理系, 中国香港, 
                {\it ``Semi-supervised Inference for Explained Variance in High-dimensional Linear Regression and Its Applications.''}, 
                2018年7月

            \item 系学术报告, 南开大学统计系, 中国, {\it ``Semi-supervised Inference for Explained Variance in High-dimensional Linear Regression and Its Applications.''}, 
                2018年7月

            \item 特邀报告, IMS 亚太会议, 新加坡, 
                {\it ``Semi-supervised Inference for Explained Variance in High-dimensional Linear Regression and Its Applications.''}, 
                2018年6月

            \item 特邀报告, 香港 EcoStat 会议, 中国香港, {\it ``Semi-supervised Inference for Explained Variance in High-dimensional Linear Regression and Its Applications.''}, 
                2018年6月

            \item 特邀报告, ICSA 2018 研讨会, 美国新布伦瑞克, {\it ``Semi-supervised Inference for Explained Variance in High-dimensional Linear Regression and Its Applications.''}, 
                2018年6月

            \item 特邀报告, 普渡大学统计研讨会, 美国, 
                {\it ``Semi-supervised Inference for Explained Variance in High-dimensional Linear Regression and Its Applications.''}, 
                2018年6月

            \item 特邀报告, 2018 杭州数据科学会议, 中国杭州, 
                {\it ``Semi-supervised Inference for Explained Variance in High-dimensional Linear Regression and Its Applications.''}, 
                2018年5月

            \item 特邀报告, 莱顿大学 Lorentz 中心, 荷兰, {\it ``Semi-supervised Inference for Explained Variance in High-dimensional Linear Regression and Its Applications.''}, 
                2018年4月

            \item 系学术报告, 哥伦比亚大学统计系, 美国, {\it ``Semi-supervised Inference for Explained Variance in High-dimensional Linear Regression and Its Applications.''}, 
                2018年4月

			\item 特邀报告, 英国剑桥, {\it ``Statistical Foundations of Uncertainty Quantification for Inverse Problems.''}, 
                2017年6月

            \item 研讨会, 宾夕法尼亚大学大数据统计方法中心, 美国, {\it ``Inference with High-dimensional Covariates and Possibly Invalid Instruments.''}, 
                2017年4月

            \item 研讨会, 坦普尔大学 Fox 商学院数据科学研究所, 美国, 
                {\it ``Inference for High-dimensional Linear Models: Fundamental Limits and Algorithms.''}, 
                2017年2月

            \item 系学术报告, 加州大学伯克利分校生物统计系, 美国, {\it ``Inference for High-dimensional Linear Models: Fundamental Limits and Algorithms.''}, 
                2017年2月

            \item 系学术报告, 罗格斯大学统计系, 美国, {\it ``Inference for High-dimensional Linear Models: Fundamental Limits and Algorithms.''}, 
                2017年2月

            \item 系学术报告, 密歇根大学统计系, 美国, {\it ``Inference for High-dimensional Linear Models: Fundamental Limits and Algorithms.''}, 
                2017年1月

            \item 系学术报告, 明尼苏达大学统计系, 美国, {\it ``Inference for High-dimensional Linear Models: Fundamental Limits and Algorithms.''}, 
                2017年1月

            \item 系学术报告, 伊利诺伊大学香槟分校统计系, 美国, {\it ``Inference for High-dimensional Linear Models: Fundamental Limits and Algorithms.''}, 
                2017年1月

            \item 系学术报告, 剑桥大学纯数学与数理统计系, 英国, {\it ``Inference for High-dimensional Linear Regression: Fundamental Limits and Algorithms.''}, 
                2017年1月

            \item 系学术报告, 加州大学圣塔芭芭拉分校统计系, 美国, 
                {\it ``Inference for High-dimensional Linear Models: Fundamental Limits and Algorithms.''}, 
                2017年1月

            \item 特邀报告, 统计数学会议, 法国弗雷瑞斯 (Fr{\'e}jus, France), {\it ``Optimal Estimation of Genetic Correlation in High-dimensional Linear Models.''}, 
                2016年12月

            \item 计量经济学午餐研讨会, 宾夕法尼亚大学经济系, 美国, {\it ``Confidence Intervals for Treatment Effects in High-dimensional Linear Models.''}, 
                2016年11月
		\end{enumerate}

		\section{\sc 其他报告}
        \begin{enumerate}
            \setlength{\itemsep}{3pt}

            \item 专题投稿报告, Joint Statistical Meetings, 美国巴尔的摩, {\it ``Optimal Estimation of Co-Heritability in High-Dimensional Linear Models''},
                2017年8月

            \item 投稿报告, Joint Statistical Meetings, 美国芝加哥, {\it ``Accuracy Assessment for High-dimensional Linear Regression''},
                2016年8月

            \item 投稿报告, Eastern North American Region (ENAR), 美国奥斯汀, {\it ``Confidence Intervals for High-Dimensional Linear Regression: Minimax Rates and Adaptivity''},
                2016年3月

            \item 海报展示, John W. Tukey Conference, 普林斯顿大学,
                {\it ``Confidence Intervals for High-Dimensional Linear Regression: Minimax Rates and Adaptivity''},
                2015年9月

            \item 投稿报告, Joint Statistical Meetings, 美国西雅图, {\it ``Distance Matrix Estimation from Noisy Observation of Low Rank Position Matrix''},
                2015年8月

            \item 投稿报告, Joint Statistical Meetings, 美国波士顿, {\it ``Instrumental Variable Approach for Mediation Analysis of Count Model''},
                2014年8月

            \item 专题投稿报告, Joint Statistical Meetings, 加拿大蒙特利尔,
                {\it `` Instrumental Variable Approach for Mediation Analysis of Zero-Inflated Count Model''},
                2013年8月

            \item 海报展示, Atlantic Causal Inference Conference (ACIC), 哈佛大学,
                {\it ``Control function Instrumental Variable Estimation of Nonlinear Causal Effect Models''},
                2013年5月
        \end{enumerate}

		\section{\sc 学术会员}

        $\bullet$ 数理统计学会 (Institute of Mathematical Statistics)

        \vspace*{-3mm}

        $\bullet$ 美国统计学会 (American Statistical Association)

        \vspace*{-3mm}

        $\bullet$ 英国皇家统计学会 (Royal Statistical Society)
        \vspace*{-3mm}

        $\bullet$ 伯努利数理统计与概率学会 (Bernoulli Society for Mathematical Statistics and Probability)
        \vspace*{-3mm}

        $\bullet$ 泛华统计协会 (International Chinese Statistical Association)

        \vspace*{-3mm}

        $\bullet$ 计量经济学会 (The Econometric Society)

        \vspace*{-2mm}
	\end{resume}
\end{document}